\subsection{Análisis y Mitigación de Riesgos}
% -----------------------------------------------------------

\subsubsection{Metodología de Evaluación}

Se identificaron y evaluaron 28 riesgos distribuidos en tres dimensiones complementarias:
\textbf{técnica} (seguridad de datos, integridad transaccional, concurrencia, disponibilidad
y calidad del software), \textbf{operativa} (errores humanos, capacitación del personal y
continuidad del negocio) y \textbf{legal} (LFPDPPP, CFF/SAT, LFPC,
responsabilidad civil y PCI DSS como riesgo evitado por alcance).

Cada riesgo fue calificado conforme a la escala que se describe a continuación. La
probabilidad refleja la frecuencia esperada de ocurrencia: \textit{Alta} ($>$50\%),
\textit{Media} (20--50\%) o \textit{Baja} ($<$20\%). El impacto pondera las consecuencias
sobre la operaci\'on, los datos personales y la exposici\'on legal: \textit{Alto}, \textit{Medio}
o \textit{Bajo}. El nivel de riesgo compuesto se determina mediante la
Tabla~\ref{tab:matriz_compuesta}.

\begin{table}[H]
\caption{Matriz de Calificación Compuesta Probabilidad $\times$ Impacto}
\label{tab:matriz_compuesta}
\centering
\small
\setlength{\tabcolsep}{6pt}
\begin{tabular}{>{\bfseries\raggedright\arraybackslash}p{2.5cm}
                >{\centering\arraybackslash}p{2.5cm}
                >{\centering\arraybackslash}p{2.5cm}
                >{\centering\arraybackslash}p{2.5cm}}
\hline
\thead{Prob. \textbackslash\ Impacto} & \thead{Bajo} & \thead{Medio} & \thead{Alto} \\
\hline
Alta  & \cellcolor{orange}ALTO    & \cellcolor{orange}ALTO     & \cellcolor{red}CR\'ITICO \\
Media & \cellcolor{altrow}MEDIO   & \cellcolor{orange}ALTO     & \cellcolor{orange}ALTO \\
Baja  & \cellcolor{green}BAJO     & \cellcolor{altrow}MEDIO    & \cellcolor{orange}ALTO \\
\hline
\end{tabular}
\end{table}

Los tipos de tratamiento aplicados son: \textbf{Mitigar} (controles t\'ecnicos u operativos
para reducir probabilidad o impacto), \textbf{Transferir} (delegar el riesgo a un tercero),
\textbf{Aceptar} (documentar y monitorear cuando el costo de mitigaci\'on supera el impacto)
y \textbf{Evitar} (eliminar la actividad que genera el riesgo).

\subsubsection{Resumen Ejecutivo de Riesgos}

La distribuci\'on por nivel de riesgo se resume en la Tabla~\ref{tab:resumen_riesgos}.
Se identificaron 28 riesgos; 27 permanecen activos y 1 queda evitado en v1.0.
Entre los activos, 4~se clasifican como cr\'iticos, 17 como altos, 5 como medios y 1 como bajo.

\begin{table}[H]
\caption{Distribución de Riesgos por Nivel}
\label{tab:resumen_riesgos}
\centering
\small
\setlength{\tabcolsep}{5pt}
\begin{tabular}{>{\bfseries\centering\arraybackslash}p{2.5cm}
                >{\centering\arraybackslash}p{1.5cm}
                >{\raggedright\arraybackslash}p{5cm}
                >{\centering\arraybackslash}p{2.5cm}}
\hline
\thead{Nivel} & \thead{\# Riesgos} & \thead{Criterio de acción} & \thead{Prioridad} \\
\hline
\rowcolor{red}
CR\'ITICO & 4 & Acci\'on inmediata antes de la prueba piloto o demostración con usuarios. & Inmediata \\ \rowcolor{orange}
ALTO     & 17 & Atender en el sprint actual o pr\'oximo. & Alta \\ \rowcolor{altrow}
MEDIO    & 5  & Planificar mitigaci\'on en los primeros 2 meses post-lanzamiento. & Media \\ \rowcolor{green}
BAJO     & 1  & Monitorear; atender si la frecuencia aumenta. & Baja \\
\rowcolor{altrow}
EVITADO  & 1  & No requiere mitigaci\'on en v1.0; el alcance se excluye por dise\~no. & N/A \\
\hline
\end{tabular}
\end{table}

\noindent Los hallazgos críticos que requieren atención obligatoria antes de la
prueba piloto o demostración con usuarios son: RT-01 y RT-02, relacionados con
concurrencia en liquidación interna; RL-01, relacionado con el cumplimiento de la
LFPDPPP; y RO-01, relacionado con la recuperación de acceso a una mesa virtual
activa. Además, se recomienda atender antes del piloto los riesgos altos RT-03 y
RT-09 por su relación directa con la seguridad de acceso a mesas virtuales.

\noindent\textit{Nota importante --- alcance de tarjetas en v1.0:} La versión~1.0 de la
plataforma no contempla integración de tarjeta ni procesamiento de datos de tarjeta.
El riesgo RL-02 queda por tanto \textbf{evitado} en esta versión; cualquier integración
posterior deberá tratarse como trabajo futuro y requerirá certificación correspondiente
o contratación de un proveedor Level~1 PCI~DSS que asuma el alcance del entorno de datos
de tarjeta (CDE).

\subsubsection{Riesgos Técnicos (RT)}

La Tabla~\ref{tab:riesgos_tecnicos} presenta los 11~riesgos t\'ecnicos identificados,
ordenados por nivel de criticidad.

\setstretch{1.2}
\setlength{\tabcolsep}{3pt}
\begin{longtable}{>{\bfseries\centering\arraybackslash}p{0.7cm}
                >{\raggedright\arraybackslash}p{4cm}
                >{\centering\arraybackslash}p{1.1cm}
                >{\centering\arraybackslash}p{1.1cm}
                >{\centering\arraybackslash}p{1.4cm}
                >{\raggedright\arraybackslash}p{6cm}
                >{\centering\arraybackslash}p{1.2cm}}
\caption{Matriz de Riesgos Técnicos} \label{tab:riesgos_tecnicos} \\
\hline
\thead{ID} & \thead{Descripción del riesgo} & \thead{Prob.} &
\thead{Impacto} & \thead{Nivel} & \thead{Estrategia de mitigación} & \thead{Trat.} \\
\hline
\endfirsthead
\caption[]{TABLA \thetable{} (Continuación)} \\
\hline
\thead{ID} & \thead{Descripción del riesgo} & \thead{Prob.} &
\thead{Impacto} & \thead{Nivel} & \thead{Estrategia de mitigación} & \thead{Trat.} \\
\hline
\endhead
\hline
\endfoot
\rowcolor{red}
RT-01 &
  Concurrencia en liquidación interna: dos comensales registran aportaciones sobre
  el mismo saldo pendiente de forma simultánea. &
  Alta & Alto & CRÍTICO &
  Usar transacciones ACID con bloqueo a nivel de fila (\texttt{SELECT FOR UPDATE}).
  Validar el saldo pendiente dentro de la transacción antes de confirmar la aportación. &
  Mitigar \\ \rowcolor{red}

RT-02 &
  Sobreliquidación en mesa compartida: aportes simultáneos superan el monto total
  pendiente. &
  Alta & Alto & CRÍTICO &
  Transacciones atómicas con bloqueo a nivel de fila (\texttt{SELECT FOR UPDATE}).
  Rechazar la operación que exceda el monto restante y retornar error al origen.
  Implementar \texttt{idempotency\_keys} por operación para evitar duplicados. &
  Mitigar \\ \rowcolor{orange}

RT-03 &
  Brecha de seguridad: credenciales de mesa virtual interceptadas en tránsito
  (ID + contraseña). &
  Media & Alto & ALTO &
  Forzar HTTPS/TLS en todas las comunicaciones. Almacenar la contraseña de mesa
  mediante hash seguro y limitar su vigencia al ciclo de vida de la mesa activa. &
  Mitigar \\ \rowcolor{orange}

RT-04 &
  Fuerza bruta sobre contraseña de mesa virtual: intentos masivos de acceso. &
  Media & Medio & ALTO &
  Aplicar rate limiting por IP y por ID de mesa, bloqueo temporal tras intentos
  fallidos y registro de eventos para auditoría. &
  Mitigar \\ \rowcolor{orange}

RT-05 &
  Inconsistencia de datos de orden: cocina visualiza una orden cancelada o
  desactualizada. &
  Media & Alto & ALTO &
  Utilizar Supabase Realtime para actualizar el tablero de cocina y definir
  transiciones de estado controladas para las órdenes. &
  Mitigar \\ \rowcolor{orange}

RT-06 &
  Degradación del flujo de liquidación interna: el comensal no puede completar
  su registro de aportación por fallo de conectividad. &
  Baja & Alto & ALTO &
  Implementar reintentos con backoff exponencial y mostrar alternativa operativa
  para registrar la aportación en caja o con apoyo del mesero. &
  Mitigar \\ \rowcolor{orange}

RT-07 &
  Gestión de sesiones deficiente: la sesión del anfitrión expira o se invalida
  durante una mesa activa, dificultando el seguimiento operativo. &
  Media & Medio & ALTO &
  Definir una duración de sesión acorde al turno operativo del restaurante.
  Permitir reingreso mediante ID de mesa y contraseña, siempre que la mesa virtual
  permanezca activa. Registrar eventos de expiración y reingreso para auditoría. &
  Mitigar \\ \rowcolor{orange}

RT-08 &
  Pérdida o corrupción de datos transaccionales por caída de base de datos durante
  una sesión activa. &
  Baja & Alto & ALTO &
  Utilizar transacciones ACID, respaldos automáticos, Write-Ahead Logging y
  mecanismos de recuperación provistos por PostgreSQL/Supabase. Definir un plan
  básico de recuperación y validar respaldos durante pruebas. &
  Mitigar \\ \rowcolor{orange}

RT-09 &
  Reutilización de enlace QR fotografiado: una persona conserva el enlace firmado
  e intenta usarlo en una sesión posterior de la misma mesa. &
  Media & Medio & ALTO &
  Validar que el acceso dependa no solo del QR firmado, sino también del estado
  activo de la mesa virtual y de la contraseña de sesión. Documentar la rotación
  de QR o versionado de mesa como mejora futura. &
  Mitigar \\ \rowcolor{orange}

RT-10 &
  Escalabilidad insuficiente en horas pico: degradación con múltiples mesas
  activas. &
  Media & Medio & ALTO &
  Realizar pruebas de carga antes del despliegue, optimizar consultas frecuentes,
  cachear lecturas del menú y separar los procesos analíticos batch para que no
  afecten la operación transaccional. &
  Mitigar \\ \rowcolor{orange}

RT-11 &
  Inyección de datos maliciosos: cantidades negativas, caracteres especiales o
  campos manipulados en órdenes y liquidaciones internas. &
  Media & Medio & ALTO &
  Validar y sanitizar todos los datos de entrada en el backend. Utilizar esquemas
  de validación, consultas parametrizadas mediante Prisma ORM, límites de cantidad
  y reglas de negocio para impedir montos negativos o valores fuera de rango. &
  Mitigar \\
\end{longtable}
\setstretch{1.5}

\subsubsection{Riesgos Operativos (RO)}

La Tabla~\ref{tab:riesgos_operativos} presenta los 9~riesgos operativos, derivados de
procesos, comportamiento humano y procedimientos de operaci\'on diaria.

\setstretch{1.2}
\setlength{\tabcolsep}{3pt}
\begin{longtable}{>{\bfseries\centering\arraybackslash}p{0.7cm}
                >{\raggedright\arraybackslash}p{4cm}
                >{\centering\arraybackslash}p{1.1cm}
                >{\centering\arraybackslash}p{1.1cm}
                >{\centering\arraybackslash}p{1.4cm}
                >{\raggedright\arraybackslash}p{6cm}
                >{\centering\arraybackslash}p{1.2cm}}
\caption{Matriz de Riesgos Operativos} \label{tab:riesgos_operativos} \\
\hline
\thead{ID} & \thead{Descripción del riesgo} & \thead{Prob.} &
\thead{Impacto} & \thead{Nivel} & \thead{Estrategia de mitigación} & \thead{Trat.} \\
\hline
\endfirsthead
\caption[]{TABLA \thetable{} (Continuación)} \\
\hline
\thead{ID} & \thead{Descripción del riesgo} & \thead{Prob.} &
\thead{Impacto} & \thead{Nivel} & \thead{Estrategia de mitigación} & \thead{Trat.} \\
\hline
\endhead
\hline
\endfoot
\rowcolor{red}
RO-01 &
  Anfitrión olvida contraseña de mesa y ningún comensal la conoce: mesa queda
  inaccesible con órdenes activas. &
  Alta & Alto & CRÍTICO &
  Implementar en v1.0 un flujo de recuperación controlado: el mesero o gerente
  podrá regenerar la contraseña de emergencia desde su rol, registrando el evento
  en auditoría y notificando a los comensales activos. &
  Mitigar \\ \rowcolor{orange}
RO-02 &
  Mesero inicializa una mesa virtual con mesas físicas incorrectas o selecciona
  una mesa que no corresponde al grupo de comensales. &
  Media & Medio & ALTO &
  Mostrar confirmación visual antes de inicializar la mesa virtual, incluyendo
  número de mesa, zona o ubicación y capacidad. Permitir corrección antes de que
  el anfitrión escanee el QR y registrar el evento de inicialización en auditoría. &
  Mitigar \\ \rowcolor{orange}
RO-03 &
  Personal sin capacitaci\'on adecuada: errores operativos desde el primer d\'ia. &
  Alta & Medio & ALTO &
  Manual de usuario por rol. Onboarding interactivo en la app (tooltips, tutorial
  guiado). Modo sandbox para pr\'actica. Superusuario de sucursal para soporte
  de primer nivel. &
  Mitigar \\ \rowcolor{orange}
RO-04 &
  Men\'u desactualizado: comensales ordenan productos descontinuados o con precios
  incorrectos. &
  Media & Medio & ALTO &
  Sincronizaci\'on peri\'odica entre el men\'u f\'isico y la plataforma. Acceso del
  gerente para actualizar en tiempo real. Alertas si un producto no ha sido
  actualizado en m\'as de 30~d\'ias. &
  Mitigar \\ \rowcolor{orange}
RO-05 &
  Chef no recibe notificaci\'on de nuevas \'ordenes en tiempo real, generando
  retrasos en cocina. &
  Media & Alto & ALTO &
  Notificaciones push o sonoras ante nuevas \'ordenes (Web Audio API).
  SLA de respuesta: confirmaci\'on de orden en $<$60~s. Alerta al gerente si una
  orden no es confirmada a tiempo. &
  Mitigar \\ \rowcolor{altrow}
RO-06 &
  Mesa compartida abandonada: comensales se van sin liquidar su parte. &
  Baja & Medio & MEDIO &
  Mesero y gerente visualizan el estatus de liquidaci\'on en tiempo real.
  Se define protocolo de validación y cierre en caja como respaldo operativo.
  La decisión de liberar o bloquear la mesa queda documentada como política del restaurante. &
  Aceptar \\ \rowcolor{altrow}
RO-07 &
  Dos personas intentan activar la misma mesa virtual como anfitri\'on
  simult\'aneamente. &
  Media & Bajo & MEDIO &
  Activaci\'on de mesa permitida una sola vez (PENDIENTE~$\rightarrow$~ACTIVA).
  El segundo intento retorna error indicando que la mesa ya fue activada. &
  Mitigar \\ \rowcolor{altrow}
RO-08 &
  Credenciales de mesa compartidas con personas externas al grupo. &
  Baja & Medio & MEDIO &
  Advertencia expl\'icita al anfitri\'on sobre no compartir la contrase\~na.
  L\'imite m\'aximo configurable de comensales por mesa.
  Mesero puede auditar el n\'umero de comensales activos. &
  Aceptar \\ \rowcolor{green}
RO-09 &
  El mesero olvida cambiar el estatus de la mesa de Liquidada a Libre. &
  Baja & Bajo & BAJO &
  Notificaci\'on autom\'atica al mesero/gerente cuando la mesa pasa a estado
  Liquidada. Transici\'on autom\'atica a Libre tras tiempo configurable de
  inactividad post-liquidaci\'on. &
  Mitigar \\
\end{longtable}
\setstretch{1.5}

\subsubsection{Riesgos Legales y Regulatorios (RL)}

El marco normativo aplicable comprende la LFPDPPP, el C\'odigo Fiscal de la
Federaci\'on (CFF), la Ley Federal de Protecci\'on al Consumidor (LFPC), el est\'andar
PCI~DSS~v4.0 y el C\'odigo Civil Federal. La Tabla~\ref{tab:riesgos_legales}
presenta los 8~riesgos legales identificados, incluyendo uno evitado en v1.0.

\setstretch{1.2}
\setlength{\tabcolsep}{3pt}
\begin{longtable}{>{\bfseries\centering\arraybackslash}p{0.7cm}
                >{\raggedright\arraybackslash}p{4cm}
                >{\centering\arraybackslash}p{1.1cm}
                >{\centering\arraybackslash}p{1.1cm}
                >{\centering\arraybackslash}p{1.4cm}
                >{\raggedright\arraybackslash}p{6cm}
                >{\centering\arraybackslash}p{1.2cm}}
\caption{Matriz de Riesgos Legales y Regulatorios} \label{tab:riesgos_legales} \\
\hline
\thead{ID} & \thead{Descripción del riesgo} & \thead{Prob.} &
\thead{Impacto} & \thead{Nivel} & \thead{Estrategia de mitigación} & \thead{Trat.} \\
\hline
\endfirsthead
\caption[]{TABLA \thetable{} (Continuación)} \\
\hline
\thead{ID} & \thead{Descripción del riesgo} & \thead{Prob.} &
\thead{Impacto} & \thead{Nivel} & \thead{Estrategia de mitigación} & \thead{Trat.} \\
\hline
\endhead
\hline
\endfoot
\rowcolor{red}
RL-01 &
  Incumplimiento LFPDPPP: recoleci\'on de datos personales sin aviso de privacidad
  ni consentimiento expl\'icito. Sanciones de hasta 320,000 d\'ias de SMGV \cite{reglamento_lfpdppp}. &
  Alta & Alto & CR\'ITICO &
  Redactar y publicar Aviso de Privacidad (Art.~16 LFPDPPP). Consentimiento expreso
  al registrarse. Nombrar un responsable de datos. Definir plazos de retenci\'on y
  procedimiento de eliminaci\'on (Art.~19). &
  Mitigar \\ \rowcolor{orange}
RL-02 &
  PCI~DSS: procesamiento de tarjeta fuera de alcance en v1.0, con exposición a sanciones
  contractuales y responsabilidad civil si se habilitara sin certificación adecuada. &
  --- & Alto & EVITADO &
  Mantener fuera del alcance de v1.0 el almacenamiento, transmisión o procesamiento
  de datos de tarjeta. Documentar cualquier integración con agregadores de pago como
  trabajo futuro mediante proveedores certificados y reducción de alcance PCI. &
  Evitar \\
RL-03 &
  Obligación fiscal SAT: los consumos registrados en la plataforma podrían requerir
  emisión de CFDI por parte del restaurante conforme al CFF Art.~29. &
  Media & Alto & ALTO &
  Consultoría fiscal para determinar responsabilidad de emisión de CFDI
  entre la cadena restaurantera y la plataforma. Documentar que la versión 1.0
  no emite facturas y que la integración con un PAC queda como trabajo futuro. &
  Mitigar \\ \rowcolor{orange}
RL-04 &
  Responsabilidad por liquidaciones internas no reconocidas o erróneas: disputa del
  comensal o del restaurante sobre el monto registrado. &
  Media & Alto & ALTO &
  Generar constancia interna de estado de cuenta, registrar sesión, mesa, productos,
  monto, usuario/comensal, fecha y hora. Publicar una política de aclaraciones
  en los Términos de Servicio y aclarar que la plataforma no procesa pagos bancarios
  reales en v1.0. &
  Mitigar \\ \rowcolor{orange}
RL-05 &
  Ausencia de T\'erminos y Condiciones de Uso: plataforma opera sin contratos
  que delimiten responsabilidades entre partes. &
  Media & Medio & ALTO &
  Redactar T\'erminos y Condiciones aceptados al acceder por primera vez.
  Especificar limitaci\'on de responsabilidad, jurisdicci\'on (M\'exico) y
  mecanismo de quejas. &
  Mitigar \\ \rowcolor{altrow}
RL-06 &
  Uso de la plataforma por menores de edad sin acompañamiento de un adulto. &
  Baja & Medio & MEDIO &
  Establecer en Términos de Uso que el acceso y autorización de consumo debe
  realizarse bajo responsabilidad del anfitrión o adulto responsable de la mesa.
  Evitar la recolección de datos personales sensibles de menores. &
  Mitigar \\ \rowcolor{orange}
RL-07 &
  Violaci\'on de derechos ARCO: usuario solicita eliminar sus datos y la
  plataforma no tiene mecanismo (Arts.~28--37 LFPDPPP). &
  Baja & Alto & ALTO &
  Funcionalidad de eliminaci\'on de cuenta y datos asociados.
  Procedimiento ARCO publicado en la app.
  Respuesta en plazo legal de 20~d\'ias h\'abiles. &
  Mitigar \\ \rowcolor{altrow}
RL-08 &
  Propiedad intelectual: uso de librer\'ias con licencias incompatibles con uso
  comercial (GPL, AGPL). &
  Baja & Medio & MEDIO &
  Auditor\'ia de licencias del stack antes de la prueba piloto o demostración con usuarios (FOSSA, licensecheck).
  Preferir MIT, Apache~2.0, BSD. Documentar todas las dependencias. &
  Mitigar \\
\end{longtable}
\setstretch{1.5}

\subsubsection{Hoja de Ruta de Mitigación}

La Tabla~\ref{tab:hoja_ruta} define el orden de atenci\'on recomendado para los riesgos
identificados, priorizando los de nivel Cr\'itico y Alto.

\setstretch{1.2}
\setlength{\tabcolsep}{4pt}
\begin{longtable}{>{\bfseries\raggedright\arraybackslash}p{3cm}
                >{\raggedright\arraybackslash}p{3cm}
                >{\raggedright\arraybackslash}p{9.5cm}}
\caption{Hoja de Ruta de Mitigación de Riesgos} \label{tab:hoja_ruta} \\
\hline
\thead{Etapa} & \thead{Riesgos a atender} & \thead{Acciones principales} \\
\hline
\endfirsthead
\caption[]{TABLA \thetable{} (Continuación)} \\
\hline
\thead{Etapa} & \thead{Riesgos a atender} & \thead{Acciones principales} \\
\hline
\endhead
\hline
\endfoot
Antes de la prueba piloto o demostración con usuarios &
  RT-01, RT-02, RT-03, RT-09, RL-01, RO-01 &
  Control de concurrencia (\texttt{SELECT FOR UPDATE}), TLS/HTTPS, QR firmado con
  HMAC-SHA256, validaci\'on de sesi\'on activa, regeneraci\'on controlada de contrase\~na
  de mesa, registro de auditor\'ia y Aviso de Privacidad conforme a la LFPDPPP. \\ \rowcolor{altrow}
Sprint~1 post-piloto (sem.\ 1--4) &
  RT-04, RT-05, RT-07, RT-11, RL-03, RL-04, RL-05 &
  Rate limiting, Supabase Realtime para tablero del Chef, gesti\'on de sesiones,
  validaci\'on de inputs, consultor\'ia fiscal SAT/CFDI, constancia interna,
  T\'erminos y Condiciones. \\
Sprint~2 post-piloto (sem.\ 5--8) &
  RT-06, RT-10, RO-02, RO-03, RO-05, RL-06, RL-07 &
  Reintentos con backoff, pruebas de carga (k6), confirmaci\'on visual de mesas,
  capacitaci\'on y onboarding, notificaciones al Chef, restricci\'on de menores,
  mecanismo ARCO. \\ \rowcolor{altrow}
Mediano plazo (mes~3--6) &
  RT-08, RO-04, RO-06, RO-07, RO-08, RO-09, RL-08 &
  Plan de recuperaci\'on ante desastres (RTO $<$ 4~hrs), gesti\'on de men\'u en
  tiempo real, pol\'iticas de retenci\'on de datos, automatizaci\'on de
  transici\'on de estatus de mesa, auditor\'ia de licencias. \\
\end{longtable}
\setstretch{1.5}

\subsubsection{Conclusiones del Análisis de Riesgos}

\noindent\textbf{Factibilidad T\'ecnica.}
La plataforma es t\'ecnicamente factible. Los riesgos t\'ecnicos identificados son manejables
mediante patrones de ingenier\'ia est\'andar (control de concurrencia, cifrado, validaci\'on
de inputs, WebSockets). El principal riesgo de implementaci\'on radica en la gesti\'on de
transacciones concurrentes (RT-01, RT-02), que debe resolverse en la capa de backend antes
de realizar una prueba piloto o demostración con usuarios mediante \texttt{SELECT FOR UPDATE} en PostgreSQL.

\noindent\textbf{Factibilidad Operativa.}
La plataforma es operativamente viable siempre que se acompa\~ne de un programa de
capacitaci\'on por rol y se definan protocolos claros para escenarios de falla. El riesgo
operativo m\'as urgente (RO-01) requiere una decisi\'on de producto antes de la
prueba piloto o demostración con usuarios:
implementar la regeneraci\'on de contrase\~na de emergencia en v1.0 y registrar el evento
en auditor\'ia.

\noindent\textbf{Factibilidad Legal.}
El cumplimiento de la LFPDPPP es no negociable y debe resolverse antes de procesar
cualquier dato personal. El riesgo RL-02 (PCI~DSS) ha sido \textbf{evitado} dentro del alcance de v1.0
mediante la decisi\'on de no integrar un agregador de pago v\'ia tarjeta ni procesar datos
de tarjeta.

\noindent\textbf{Veredicto.}
El proyecto es factible en sus tres dimensiones, condicionado a la resoluci\'on de los
4~riesgos cr\'iticos antes de la prueba piloto o demostración con usuarios. Con las mitigaciones propuestas implementadas,
el perfil de riesgo residual de la plataforma es aceptable para una prueba controlada,
demostraci\'on acad\'emica y eventual despliegue piloto.

% ============================================================
%  CAPÍTULO 4 — METODOLOGÍA Y ANÁLISIS
% ============================================================
\newpage
